% Chapter 2: What Trapdoor Computing Is Not
% DEFINES: ch:what-it-is-not; sec:privacy-source, sec:not-oram, sec:not-fhe,
%   sec:not-sse, sec:comparison-guardrail, sec:ch2-notes; tab:not
% REFERENCES (forward, part labels only): part:abstraction,
%   part:confidentiality, part:constructions

\chapter{What Trapdoor Computing Is Not}
\label{ch:what-it-is-not}

The previous chapter described a service that computes on encodings it
cannot read. That description is close enough to several well-known ideas,
oblivious RAM, homomorphic encryption, garbled circuits, searchable
encryption, that a reader could reasonably assume trapdoor computing is one
of them, or a repackaging of them. It is not, and the differences are not
cosmetic. This chapter draws the lines, and in doing so fixes the
guardrail that the rest of the book leans on: a precise statement of where
the privacy comes from, against which every later claim can be checked.

\section{Where the Privacy Comes From}
\label{sec:privacy-source}

The single most important sentence in this book is about a negative.
Trapdoor computing's confidentiality comes from \emph{one-way hashing and
uniform representation}, and not from hiding the pattern of access, and not
from a reduction to a hard problem. The untrusted machine is allowed to see
which entries of its table it touches, in what order, and how often. What
it is denied is the \emph{meaning} of those entries: a token is the image
of a value under a keyed one-way encoding, and many values are spread
across uniform-looking tokens, so the access it observes is access to
opaque, frequency-flattened strings. Privacy is a property of the
\emph{representation}, not of the protocol that moves it around.

This is worth stating sharply because the neighboring approaches each hide
something else, at a different cost, and conflating them imports the wrong
intuitions. The rest of the chapter takes them one at a time, fairly: each
solves a real problem that trapdoor computing does not, and the point is
the difference in what is hidden, never that one is simply better.

\section{Not Oblivious RAM}
\label{sec:not-oram}

Oblivious RAM hides the \emph{access pattern}: it guarantees that the
sequence of memory locations a computation touches is independent of the
data and the query, so an observer who sees the addresses learns nothing
from them. It pays for this with polylogarithmic overhead per access and a
stateful, interactive protocol. Trapdoor computing makes the opposite
trade. It does \emph{not} hide the access pattern, the untrusted machine
plainly sees which table entries are read, and it spends nothing to conceal
them. What it hides is what those entries mean. An ORAM and a cipher map,
pointed at the same encrypted store, defend against different observers: the
ORAM defeats one who reads addresses, the cipher map defeats one who reads
contents. Prior work in this framework occasionally drifted toward
access-pattern language, and the result was confusion; this book does not,
and a claim that ever rests on access-pattern indistinguishability has
wandered off the framework.

\section{Not Homomorphic Encryption, Not Garbled Circuits}
\label{sec:not-fhe}

Fully homomorphic encryption computes \emph{exactly} on ciphertexts: any
function of the plaintexts can be evaluated under encryption, with a result
that decrypts to the true value. Its guarantee is semantic security of the
ciphertexts, and its cost, despite a decade of progress, remains large in
both ciphertext size and per-operation compute. A cipher map is exact in
neither sense and asks for neither cost. It \emph{approximates}: the
decoded result equals the latent value only up to a bounded error rate, and
that error is not a defect to be apologized for but, as the previous
chapter noted, a source of deniability. Where homomorphic encryption buys
exactness at high cost, trapdoor computing buys cheapness and a measured
leak by giving exactness up on purpose.

Garbled circuits sit closer in spirit, since they too let one party
evaluate a function blindly, but they are built for a different shape of
problem. A garbled circuit hides a function's inputs for a \emph{single}
evaluation, paid for with cryptographic work at every gate, and is consumed
in the using; a fresh garbling is needed for the next input. A cipher map
is a \emph{static, total lookup table}, built once and reused across as
many queries as the deployment runs, with no per-query cryptography beyond
the encode and decode at the trusted end. The garbled circuit is a
one-shot interactive object; the cipher map is a durable artifact you can
ship to an untrusted server and query indefinitely.

\section{Not (Just) Searchable or Property-Preserving Encryption}
\label{sec:not-sse}

The nearest neighbor is encrypted search. Searchable symmetric encryption
and property-preserving encryption let a server answer queries over
encrypted data, and like trapdoor computing they accept that the server
learns \emph{something}. The difference is in how that something is
governed. In the searchable-encryption tradition, the leakage is a
structural side effect, an encrypted index whose shape, or an
order-preserving ciphertext whose comparisons, the server can read, and a
long line of leakage-abuse attacks has shown how much those side effects
give away in practice. Here, leakage is not a side effect to be bounded
after the fact but the \emph{design variable}: it is the entropy ratio of
the previous chapter, a measured quantity the constructions are built to
drive down.\footnote{The measure and the constructions that move it are
developed in \Cref{part:confidentiality} and realized in
\Cref{part:constructions}.}

Frequency hiding sharpens the contrast. Frequency-smoothing encryption
defends a deterministic scheme against frequency analysis by giving common
values several ciphertexts, so the observable distribution flattens; this
is the right idea, and it is the same idea trapdoor computing uses, more
frequent values get more representations, with the allocation $K(x)$ grown
in proportion to a value's frequency $D(x)$. The difference is mechanism.
In frequency-smoothing encryption the flattening is enacted by an online
protocol that injects or rewrites ciphertexts as data arrives; here it is a
static property of a total function, present in the table the untrusted
machine holds, with no protocol to run and no traffic to shape. The lever
is the same; the place it is pulled is not.

\section{The Comparison, and the Guardrail}
\label{sec:comparison-guardrail}

\Cref{tab:not} collects the distinctions. Read down the ``hides'' column:
each approach conceals a different thing, and the costs differ accordingly.
Trapdoor computing's row is the one that hides \emph{meaning} via a total
map, pays in the space it takes to separate signal from noise and to encode
values, and reports its residual leak as a number rather than a promise.

\begin{table}[t]
\centering
\caption{Trapdoor computing against neighboring approaches. Each hides
something different; none is uniformly stronger.}
\label{tab:not}
\begin{tabular}{@{}p{2.4cm}p{2.8cm}p{3.0cm}p{3.4cm}@{}}
\toprule
Approach & Hides & Cost & Different here \\
\midrule
Oblivious RAM & the access pattern & polylog overhead per access; interactive & access is not hidden, meaning is \\
Homomorphic encryption & the plaintext & large ciphertexts and compute & we approximate, measurably; FHE is exact \\
Garbled circuits & inputs, one evaluation & per-gate cryptography, one-shot & one static table, reused across queries \\
Searchable / property-preserving encryption & data and queries, partially & structured index or order leakage & leakage is the design variable, measured \\
\midrule
Trapdoor computing & meaning, via a total map & space, $-\log_2 \varepsilon + \mu$ bits per element & privacy from hash and uniformity, leak is the entropy ratio \\
\bottomrule
\end{tabular}
\end{table}

This table is also the book's guardrail. The framework's privacy rests on
one-way hashing and uniform representation, and on nothing else; in
particular this book imports no notion of access-pattern
indistinguishability, no differential-privacy mechanism, no
simulation-based or game-based security definition. Those are the right
tools for the problems in the table's other rows, and reaching for them
here is the standing way the framework gets misstated. The reader, and the
author, should treat any later argument that quietly depends on one of them
as a drift to be corrected, not a result to be kept.

\section{Notes and Provenance}
\label{sec:ch2-notes}

Oblivious RAM is due to Goldreich and Ostrovsky~\cite{goldreich1996software},
with Path ORAM~\cite{stefanov2013path} the now-standard efficient
construction; fully homomorphic encryption to Gentry~\cite{gentry2009fully};
garbled circuits to Yao~\cite{yao1986generate}. Searchable symmetric
encryption begins with Song, Wagner, and Perrig~\cite{song2000practical} and
the definitions of Curtmola et
al.~\cite{curtmola2006searchable}; the leakage-abuse line that motivates
treating leakage as a first-class quantity runs through Islam et
al.~\cite{islam2012access}, Naveed et al.~\cite{naveed2015inference}, and
Cash et al.~\cite{cash2015leakage}; frequency-smoothing encryption is
Lacharit\'e and Paterson~\cite{lacharite2018frequency}. The classification
of which sources in this framework are authentic and which were expanded by
later drafting is maintained in the project's ecosystem triage; the
correction of the access-pattern drift named in
\Cref{sec:not-oram} is one of its recorded items.
